\section{Discussion}
\label{sec:disc}

\subsection{Why Strong Models Fail Silently}
\label{sec:disc-why}

On synthetic trap tasks, the corrected SFR is high (86--100\%) because these tasks are designed so that code runs but the answer is wrong. On real UCI data, SFR drops to 4\%, and on DSBench to 19\%. A null-model analysis ($R^2 = 0.493$, descriptive only) shows part of this variation is a mechanical consequence of execution-success differences: stronger models crash less, so their remaining failures are more likely silent. However, SFR and execution-success share the loud-failure count, so this correlation is partially mechanical by construction and cannot be decomposed into ``mechanical'' vs.\ ``real'' components. The practical implication is that improving code-generation ability alone will shift the failure distribution toward silent failures, but silent failures are most prevalent on tasks where code execution succeeds but reasoning fails.

\subsection{Implications for Agent Deployment}
\label{sec:disc-impl}

The zero recovery rate (\Cref{sec:res-recovery}) has direct deployment consequences. An agent that never self-corrects silent failures will propagate wrong answers through multi-step pipelines, and the errors compound. Practitioners cannot rely on the agent to ``notice'' it is wrong; external verification is mandatory. However, our verifier ablation (\Cref{sec:res-verifier}) shows that simple rule-based verification is a double-edged sword: it helps weak models but not strong ones. Effective verification likely requires model-specific calibration or learned verifiers.

\subsection{Negative Control: Oracle Repair Revisited}
\label{sec:disc-negative-control}

To test whether the oracle repair rate (\Cref{sec:causal}) genuinely reflects step-specific causal attribution, we ran a negative control on 56 failed tasks. The \emph{no-op} intervention runs the ground-truth code independently in a fresh sandbox, without any step replacement. As \Cref{tab:negative-control} shows, the two rates are identical (64.3\%), and the conservative true-positive rate is 0.0\%. This confirms that the oracle repair rate measures whether the ground-truth code produces the correct answer, \emph{not} whether replacing the originating step has a step-specific causal effect. We therefore report oracle repair as an upper bound on reparability rather than as evidence of root-cause localisation.

\begin{table}[t]
\centering
\caption{Negative control analysis. The no-op intervention (running GT code independently) produces identical results to the oracle intervention, confirming that oracle repair measures GT code correctness, not step-specific causal effect.}
\label{tab:negative-control}
\small
\begin{tabular}{@{}lcc@{}}
\toprule
Intervention & Tasks repaired & Rate (\%) \\
\midrule
Oracle (replace originating step) & 36/56 & 64.3 \\
No-op (run GT code independently) & 36/56 & 64.3 \\
\midrule
Conservative TP (oracle $\wedge$ $\neg$noop) & 0/56 & 0.0 \\
\bottomrule
\end{tabular}
\end{table}


A detailed discussion of limitations and ethical considerations is provided in \Cref{sec:limits-ethics}.
